% Source: https://artofproblemsolving.com/community/c3h2681798_is_there_a_trick
% Stable identity: aops-topic-2681798-q01
\emph{The trick: simplify $g$ first.} $g(x)=\dfrac{4(x+4)}{(x+1)(x+4)}=\dfrac4{x+1}$ for $x\ne-4,-1$ ($g$ is undefined at both).

\emph{$f(g(x))$.} Put $u=\dfrac4{x+1}$; we need $u\ne5$, i.e. $x\ne-\dfrac15$. Multiplying the numerator and denominator of $f(u)$ by $(x+1)^2$:
\[f(g(x))=\dfrac{16-12(x+1)-3(x+1)^2}{(x+1)\left(4-5(x+1)\right)}=\dfrac{-3x^2-18x+1}{-(x+1)(5x+1)}=\boxed{\dfrac{3x^2+18x-1}{(x+1)(5x+1)}},\]
defined for $x\ne-4,-1,-\dfrac15$.

\emph{$g(f(x))$.} We need $x\ne5$ and $f(x)\ne-4,-1$. Then $g(f(x))=\dfrac4{f(x)+1}$ and
\[f(x)+1=\dfrac{x^2-3x-3+x-5}{x-5}=\dfrac{(x-4)(x+2)}{x-5},\qquad\text{so}\qquad g(f(x))=\boxed{\dfrac{4(x-5)}{(x-4)(x+2)}}.\]
Here $f(x)=-1$ at $x=4$ and $x=-2$, and $f(x)=-4$ when $x^2-3x-3=-4x+20$, i.e. $x^2+x-23=0$, $x=\dfrac{-1\pm\sqrt{93}}2$. So $g(f(x))$ is defined for $x\ne5,4,-2,\dfrac{-1\pm\sqrt{93}}2$. (At $x=\dfrac{-1\pm\sqrt{93}}2$ the simplified formula has a value, but the composition does not, because $g$ is undefined at $-4$.)
